\documentclass[10pt,a4paper]{amsart}
\usepackage[margin=19mm]{geometry}
\usepackage{amsmath,amssymb,mathtools}
\usepackage[scr=rsfs,cal=euler]{mathalfa}
\usepackage{xfrac}
\usepackage[unicode,hidelinks,bookmarks=false]{hyperref}
\newcommand{\ie}{i.e.\ }
\newcommand{\cf}{cf.\ }
\newcommand{\resp}{respectively\ }
\newcommand{\Subsection}{\subsection}
\providecommand{\nobreakdash}{\nobreak\hbox{-}}
\setlength{\emergencystretch}{2em}
\multlinegap0pt
\usepackage{amssymb}
\usepackage{tikz}
\newtheorem{theorem}{Theorem}
\newcommand\Fcal{{\mathcal F}}
\newcommand\Kcal{{\mathcal K}}
\DeclareMathOperator{\Sal}{Sal}
\newcommand\Tcal{{\mathcal T}}
\newcommand\Ztwo{{\mathbb Z}/2{\mathbb Z}}
\newcommand\dd{{\bf d}}
\newcommand\ff{{\bf f}}
\newcommand\uu{{\bf u}}
\newcommand\vv{{\bf v}}
\title{Filtrations of tope spaces of oriented matroids}
\author{Kris Shaw, Chi Ho Yuen}
\date{Source: arXiv:2501.11295v2 (CC BY 4.0)}
\begin{document}
\maketitle

\noindent\emph{Source and licence.} Filtrations of tope spaces of oriented matroids. Version arXiv:2501.11295v2 (CC BY 4.0), distributed by arXiv under the Creative Commons Attribution 4.0 International licence, http://creativecommons.org/licenses/by/4.0/, as recorded in the arXiv licence metadata for this exact version and shown on the abstract page https://arxiv.org/abs/2501.11295v2. Full licence notice: \texttt{licenses/arxiv-2501.11295-CC-BY-4.0.txt}.\par\smallskip

\noindent\textit{Editorial scope.} Counted unit: the article's theorem thm:Kalinin\_prefix (source0.tex lines 716--719). The article's complete original proof of it is delivered with it (source0.tex lines 721--791, one \texttt{proof} environment of 71 lines). No mathematics is written, repaired or re-derived: the statement and the proof are the article's own lines, in the article's own notation, and no retained line is substituted or re-flowed.  No further source-local context is needed: every cross-reference of every retained line resolves inside the two delivered spans.  Nothing was omitted from any retained span, so the package carries no omission record.  No retained line contains a citation, so the wrapper carries no bibliography and no bibliography entry.  No retained line contains a figure environment, an image-inclusion command or drawing code, so no image is delivered and the package declares no asset dependency.  Editorial formatting only, in addition: an AMS \texttt{amsart} preamble that restates the article's own declarations and macros used by the retained lines, namely the environments \texttt{theorem} on separate counters, together with the standard packages those lines need.  The retained environments print in this document as Theorem 1 in order of appearance, whereas the article prints the same items with the numbering of its own full document.  The complete native source of the article as distributed in the arXiv e-print for this exact version is bundled unchanged as \texttt{sources/171828/source0.tex}.
\begin{theorem} \label{thm:Kalinin_prefix}
Let $\Fcal:\emptyset=F_0 \subsetneq  F_1 \subsetneq \ldots \subsetneq F_d=E$ be a complete flag of flats and let $T_{\vv}\in\Tcal_\Fcal$.
Then $\gamma_{\Fcal,\vv,p}\in\Kcal_p(M)$.
\end{theorem}

\begin{proof}
For brevity of the notation, we assume (by reorientation) $\vv={\bf 0}$ whenever only one specific prefix chain is being considered.

For the case $p=1$, we have $\gamma_{\Fcal, {\bf 0},1}=T_{{\bf 0}}+T_{\dd_1}$, while $Z(T_{\bf 0}\setminus F_1,T_{\bf 0})$ is a $1$-dimensional cell whose boundary is precisely the two vertices $Z(T_{\bf 0},T_{\bf 0})$ and 
$Z(T_{\dd_1},T_{\bf 0})$ .
So we can take $\beta_1$ to be $Z(T_{\bf 0}\setminus F_1,T_{\bf 0})$.

Now suppose the statement is true for $p$. More precisely, for every $\gamma_{\Fcal,\vv,p}$ we suppose for $i \leq p$ that 
there exists $ \beta_i \in C_i(\Sal_M; \Ztwo)$ such that 
$\partial\beta_1= i_{\ast} \gamma$ and $\partial\beta_i=\beta_{i-1}+\overline{\beta_{i-1}}$ for $ 2 \leq i \leq p$. 
We further assume that we can choose 
\begin{equation}\label{eq:beta_p}
    \beta_{p}  = \sum_{{\bf u}\in U_{\Fcal,\vv,p-1}} Z(T_{\vv}\setminus F_{p},T_{\bf u}).
\end{equation}

Notice that this holds in the case $p = 1$ above.

We now pass to the case of $p+1$, and prove that $\gamma_{\Fcal,{\bf 0},p+1}$
is in $\Kcal_{p+1}(M)$ by finding $i$-chains $\beta_i$ for $i = 1, \dots, p+1$. Moreover, we have 
$$\beta_{p+1} = \sum_{{\bf u}\in U_{\Fcal, {\bf 0}, p}} Z(T_{\bf 0}\setminus F_{p+1},T_{\bf u}).$$
 
We start with noting $$\gamma_{\Fcal, {\bf 0},p+1}=\gamma_{\Fcal, {\bf 0},p}+\gamma_{\Fcal, \dd_{p+1},p}.$$ 
By induction, for $i = 1, \ldots, p$, 
there are $i$-chains $\beta_i'$ and $\beta_i''$ certifying that $\gamma_{\Fcal, {\bf 0},p}$ and $\gamma_{\Fcal, \dd_{p+1},p} $ are in $\Kcal_p$.
For $i \leq p$, we can take $\beta_i = \beta_i' + \beta_i''$. Using the induction assumption, it is easily verified that $\partial \beta_i = \beta_{i-1} +\overline{\beta_{i-1}}$ for $i = 1, \ldots, p$. 
It remains to show that $\partial \beta_{p+1} = \beta_p + \overline{\beta_{p}}$, where $\beta_p = \beta_p' + \beta_p''$.

We spend some time simplifying the conjugation of $\beta'_p$. 
Observe that
$$\overline{Z(T_{\bf 0}\setminus F_p,T_{\bf u})}
=Z(T_{\bf 0}\setminus F_p,-T_{\bf u})
=Z(T_{\bf 0}\setminus F_p,T_{\ff_d + \bf u})
=Z(T_{\bf 0}\setminus F_p,T_{\ff_p+ {\bf u}}),$$ 
since $(T_{\bf 0}\setminus F_p) \circ T_{\ff_d + \bf u} = (T_{\bf 0}\setminus F_p) \circ T_{\ff_p+ {\bf u}}$.
Hence we have
 $$\overline{\beta'_p} = \sum_{{\bf u}\in U_{\Fcal,\ff_p,p-1}} Z(T_{\bf 0}\setminus F_p,T_{\bf u}).$$

By the same argument, we also have 
$\gamma_{\Fcal, \dd_{p+1},p}\in\Kcal_p(M)$.
More precisely, $\gamma_{\Fcal, \dd_{p+1},p}$ can be obtained from $\gamma_{\Fcal, {\bf 0},p}$ by replacing every $T_{\bf v}$ by $T_{{\bf v}+\dd_{p+1}}$ in the sum.
By induction, when constructing $\beta''_{p'}, p'<p$ for $\gamma_{\Fcal, \dd_{p+1},p}$, one can simply apply the analogous replacement of $Z(T_{\bf 0}\setminus F_{p'},T_\vv)$ by $Z(T_{\dd_{p+1}}\setminus F_{p'},T_{{\bf v}+\dd_{p+1}})$.

In particular, we can choose the corresponding $p$-chain to be 
$$\beta_p'' = \sum_{{\bf u}\in U_{\Fcal,\dd_{p+1}, p-1}} Z(T_{\dd_{p+1}} \setminus F_p,T_{\bf u}).$$

Nevertheless, we note that $Z(T_{\dd_{p+1}} \setminus F_p,T_{\uu+\dd_{p+1}})=Z(T_{\dd_{p+1}} \setminus F_p,T_{\bf u})$ for ${\bf u}\in U_{\Fcal,{\bf 0},p-1}$, as $$(T_{\dd_{p+1}}\setminus F_p)\circ T_{{\bf u}+\dd_{p+1}}=(T_{\dd_{p+1}}\setminus F_p)\circ T_{\bf u},$$ so the sum can actually be rewritten as
\begin{equation} \label{eq:2nd_chain}
  \beta_p'' =  \sum_{{\bf u}\in U_{\Fcal,{\bf 0}, p-1}} Z(T_{\dd_{p+1}} \setminus F_p,T_{\bf u}).
\end{equation}
Analogous to the above, 
we have $$\overline{\beta''_p} = \sum_{{\bf u}\in U_{\Fcal,\ff_p, p-1}} Z(T_{\dd_{p+1}} \setminus F_p,T_{\bf u}).$$

Hence, it remains to show that the boundary of
$\beta_{p+1}$ is equal to 
\begin{equation} \label{eq:target_chain}
    \sum_{{\bf u}\in U_{\Fcal,{\bf 0}, p}}\big[Z(T_{\bf 0}\setminus F_p,T_{\bf u})+Z(T_{\dd_{p+1}} \setminus F_p,T_{\bf u})\big],
\end{equation}
here we use the simple fact that $U_{\Fcal,{\bf 0}, p}=U_{\Fcal,{\bf 0}, p-1}\sqcup U_{\Fcal,\ff_p, p-1}$.

Every $p$-cell on the boundary of $\beta_{p+1}$ is of the form $Z(T\setminus F,T_{\bf u})$ for some flat $F\subset F_{p+1}$ of rank $p$, some tope $T$ that agrees with $T_{\bf 0}$ outside of $F_{p+1}$, and some ${\bf u}\in U_{\Fcal, {\bf 0}, p}$.

Case I: $F\neq F_p$.\\
There exists some $1\leq k\leq p$ such that $F\cap(F_k\setminus F_{k-1})=\emptyset$, for otherwise we would have $F_p = \langle F\cap(F_i\setminus F_{i-1}): i=1,\ldots,p\rangle\subsetneq F$, which is impossible by considering rank.
For any ${\bf u}\in U_{\Fcal,{\bf v}_0, p}$, we have $Z(T\setminus F,T_{\bf u})=Z(T\setminus F,T_{{\bf u}+\dd_k})$ because $(T\setminus F)\circ T_{\bf u}=(T\setminus F)\circ T_{{\bf u}+\dd_k}$.
By pairing up the elements of $U_{\Fcal,{\bf 0}, p}$ into $\{{\bf u}, {\bf u}+\dd_k\}$'s, we can see that the boundary does not contain cells of the form $Z(T\setminus F,T_{\bf u})$ for $F\neq F_p$. 

Case II: $F=F_p$.\\
In such a case, a tope $T$ restricted to $F_{p+1}\setminus F_p$ must be equal to that of $T_{\bf 0}$ or $T_{\dd_{p+1}}$, for otherwise, $T|_{F_{p+1}}/F_p$ would be a third tope of the rank 1 matroid $M|_{F_{p+1}}/F_p$.
For every ${\bf u}\in U_{\Fcal,{\bf 0}, p}$, the cell $Z(T_{\bf 0}\setminus F_{p+1},T_{\bf u})$ is the unique cell in $\beta_{p+1}$ whose boundary contains $Z(T_{\bf 0}\setminus F_p,T_{\bf u})$; it is also the unique cell in $\beta_{p+1}$ which contains $Z(T_{\dd_{p+1}}\setminus F_p,T_{\bf u})$ in its boundary: for any other ${\bf u}'\in U_{\Fcal,{\bf 0}, p}$, we have $(T_{\bf 0}\setminus F_p)\circ T_{\bf u}\neq (T_{\bf 0}\setminus F_p)\circ T_{{\bf u}'}$, so $Z(T_{\bf 0}\setminus F_p,T_{\bf u})\neq Z(T_{\bf 0}\setminus F_p,T_{{\bf u}'})$, and an analogous argument for $Z(T_{\dd_{p+1}}\setminus F_p,T_{\bf u})$.
Hence the two cells $Z(T_{\bf 0}\setminus F_p,T_{\bf u}),Z(T_{\dd_{p+1}}\setminus F_p,T_{\bf u})$ appear in the boundary exactly once each, as claimed in (\ref{eq:target_chain}).
\end{proof}

\end{document}
